\section{A cycle with two pendant leaves}\label{sec:leaves}

This section proves Theorem~\ref{thm:leaves}, using the same branch count
and the equality case of the positive argument.

\begin{lemma}[Restriction does not create branches]\label{lem:restriction}
Fix some inputs of a normalized circuit to constants. Propagate constants,
suppress resulting projections and repeated-operand gates, and remove
dead gates. If the resulting function is nonconstant, its normalized
circuit has no more branching wires than the original.
\end{lemma}
\begin{proof}
Deleting a constant wire or a dead gate removes vertices and arcs and
cannot increase an outdegree. The only other graph operation suppresses
a gate $q$ whose remaining value is its single parent $p$. Redirect all
uses of $q$ to $p$, including the selected output if necessary. Write
$d$ for the number of consumers of $p$ other than $q$, and $t$ for the
number of consumers of $q$. The new outdegree of $p$ is at most $d+t$;
coincident uses may reduce it. If neither $p$ nor $q$ previously branched,
then $d=0$ and $t\le1$, so the merged wire does not branch. Otherwise the
one merged wire contributes at most the previous two branch indicators.
All other outdegrees are unchanged or decrease. These operations implement
the local normalization after restriction, proving the claim by induction
on their number. Arbitrary resynthesis of the restricted function is not
part of the operation asserted here.
\end{proof}

\begin{corollary}[An induced core]\label{cor:core}
Let $G$ have $n$ vertices and no isolated vertex. Suppose an induced
subgraph $H$ is nonempty, has minimum degree at least two, and has no
four-cycle. Then
every circuit for $F_G$ has at least $n-1+\tau(H)$ AND/OR gates.
\end{corollary}
\begin{proof}
Normalize the circuit for $G$. Set every input outside $H$ to zero and
use the local simplification of Lemma~\ref{lem:restriction}. The remaining
function is exactly $F_H$. The proof of Theorem~\ref{thm:binary} requires
at least $\tau(H)$ branches in that circuit, so at least as many were
present in the original normalized DAG. Apply Lemma~\ref{lem:budget}
there, with all $n$ original inputs essential.
\end{proof}

\begin{proof}[Proof of Theorem~\ref{thm:leaves}]
The graph has $n=2k+2$ vertices and $2k+2$ edges. Its induced cycle has
cover number $k$. Thus every normalized circuit satisfies
\begin{equation}\label{eq:leafbudget}
                   g\ge3k+1+t,
\end{equation}
where $t$ counts NOT gates.

The only covers of the cycle of size $k$ are its two alternating
bipartition classes. Indeed its complement would be an independent set
of size $k$; around the cycle every gap between consecutive selected
vertices must have exactly one omitted vertex. If the two roots have
the same parity, that class covers the enlarged graph, so $\tau(G)=k$.
If they have opposite parity, no cover of the enlarged graph has size
$k$: its cycle part must already have size $k$, leaving neither leaf
available, and would therefore need both roots. One alternating cover
together with the other root is a cover of size $k+1$.

For completeness, the cover construction assigns each edge to one of
its endpoints in a minimum cover $S$. Every $v\in S$ has a neighbor
outside $S$, since otherwise it could be removed from $S$, so it receives
at least one edge regardless of how edges within $S$ are assigned.
A star with $d$ assigned edges costs $d$ gates: $d-1$ OR gates on its
neighbors and one AND with its center. Joining the $|S|$ stars uses
$|S|-1$ further gates, for a total $|E|+|S|-1$.
This gives upper bounds $3k+1$ and $3k+2$ in the two cases.
Equation~\eqref{eq:leafbudget} settles the same-parity case.

Now suppose the roots have opposite parity, and normalize a circuit of
at most $3k+1$ gates. Equation~\eqref{eq:leafbudget} forces exactly
$3k+1$ gates and $t=0$, so this normalized circuit is positive.
Let $S$ be the eligible \emph{cycle} vertices. Every cycle edge has
endpoints of degree at least two. Lemma~\ref{lem:carrier} therefore
shows that $S$ covers the cycle, and that its first branching carriers
are distinct with singleton supports. There are at least $k$ of them.
Equality in the incidence budget permits only $k$ branches altogether.
Consequently $|S|=k$, $S$ is one alternating class, and these are all the
branches. Each has exactly two consumers. No leaf is eligible, because
each branch carries only its selected cycle vertex.

Exactly one root belongs to $S$; call it $v$. As in the proof of
Lemma~\ref{lem:slots}, only the two exits of $b(v)$ can
lead to gates consuming $v$, so at most two gates consume it. Its three
incident edges each have a meeting AND gate consuming $v$.
The meeting gate exists for any edge, including a leaf edge: take the
first wire true on that pair and false on both singletons, as in the
first paragraph of the proof of Lemma~\ref{lem:carrier}.
These three gates are distinct. If two coincided, their parent not
carrying $v$ would carry both other endpoints. Among any two of the
three neighbors, at least one is an ineligible cycle vertex. Every
carrier of that vertex is a cut carrier, so Lemma~\ref{lem:carrier}
makes its support a singleton, a contradiction. Thus three gates
consume $v$, which is impossible. The opposite-parity case requires
$3k+2$ gates.
\end{proof}

The eight-cycle in Theorem~\ref{thm:leaves} uses opposite classes: vertices $0$ and $3$
have different parity and are not adjacent. An explicit fourteen-gate
witness, returning wire $23$ from primary wires $0,\ldots,9$, is
\begin{center}
\begin{minipage}{0.9\linewidth}
\centering
\begin{tabular}{rl@{\qquad}rl}
\toprule
Wire&Operation&Wire&Operation\\
\midrule
10&$1\OR7$&11&$10\OR8$\\
12&$0\AND11$&13&$1\OR3$\\
14&$2\AND13$&15&$4\OR9$\\
16&$3\AND15$&17&$4\OR6$\\
18&$5\AND17$&19&$6\AND7$\\
20&$12\OR14$&21&$16\OR18$\\
22&$20\OR21$&23&$22\OR19$\\
\bottomrule
\end{tabular}
\end{minipage}
\end{center}
Its five stars have centers $0,2,3,5,7$ and partition the ten edge terms.
Thus distributivity proves correctness and the listed rows establish the
gate count directly.
